% ARCHIVED 2026-10-13 from _c_framework.tex (Section 3.5). Cut from the paper: round trips (M, N), their
% reconstruction errors, the round-trip bounds, and the perfect-reconstruction theorems feed no application result.
% Depends on archive/c_characterizing_P.tex (overlap, ambiguity). Kept for possible reintroduction.

\subsection{Quantifying Explanation Transformation Preservation}
\label{sec:Framework:Preservation}

With operators capable of moving explanations from one space to another, we now consider the quality (or lack thereof) of a transferred explanation. This is a difficult question to answer, as the best way to measure this quality would be to compare a transferred explanation to the `correct' explanation that should have been produced, but obviously, we emphasize a practical use-case where no such explanation exists. To this end, we propose an intrinsic evaluation criteria based around the concept of cyclic transfers, i.e., moving an explanation from one space and back again.

\definition[\textit{Cyclic Explanation Transfers}]
\label{def:CyclicTransfers}
{
Consider $a \in \mathcal{A}$, $b \in \mathcal{B}$, $\tau: \mathcal{A} \mapsto \mathcal{B}$, and $\psi: \mathcal{B} \mapsto \mathcal{A}$. We define the pull--push (pull, then push) cyclic operator $\psi_\tau: \mathcal{A} \mapsto \mathcal{A}$ as $\psi_\tau(a) \coloneq \psi(\tau(a)) = M a$ with $M = P^\top D^\dagger P$, and the push--pull (push, then pull) operator $\tau_\psi: \mathcal{B} \mapsto \mathcal{B}$ as $\tau_\psi(b) \coloneq \tau(\psi(b)) = N b$ with $N = P P^\top D^\dagger$.
}

The entry $M_{k,j} = \sum_i P_{i,j} P_{i,k} / D_{i,i}$ is the share of feature $j$'s attribution that a round trip through the elements lands on feature $k$; in the probabilistic reading above, it is the probability that the problem-space mistakes feature $j$ for feature $k$. Likewise, column $i$ of $N$ is where element $\bar{z}_i$'s attribution lands after a round trip through the features.

\definition[\textit{Cyclic Reconstruction Error}]
\label{def:CyclicTransferError}
{
Consider $a \in \mathcal{A}$, $b \in \mathcal{B}$, $\psi_\tau: \mathcal{A} \mapsto \mathcal{A}$, and $\tau_\psi: \mathcal{B} \mapsto \mathcal{B}$. We define the pull--push error as $\delta_{\psi_\tau} = \|\psi_\tau(a) - a\|_1$ and the push--pull error as $\delta_{\tau_\psi} = \|\tau_\psi(b) - b\|_1$.
}

We measure the errors in the $\ell_1$ norm because it counts attribution mass: an error of $0.3\,\|a\|_1$ means that $30\%$ of the attribution ended up somewhere else. Definition~\ref{def:ExplanationProvenance} provides no guarantee that either error is small. Nonzero cyclic error suggests that at least one of the operations muddies the transfer; zero cyclic error, a \textit{perfect} cyclic transfer, suggests that both preserve the explanation. The structure of $M$ and $N$ determines both errors.

\proposition[\textit{Round Trips}]{
\label{prop:RoundTrips}
The pull--push matrix $M$ is symmetric, nonnegative, and doubly stochastic, its eigenvalues lie in $[0, 1]$, and its diagonal is $M_{j,j} = 1 - O_j$. The push--pull matrix $N$ is nonnegative, its columns sum to one for every element with $D_{i,i} > 0$, and its diagonal $N_{i,i} = \sum_j P_{i,j}^2 / D_{i,i}$ is at most $1$, with equality if and only if every feature that $\bar{z}_i$ influences draws on $\bar{z}_i$ alone.
\textit{Proof.} $M_{k,j} = \sum_i P_{i,j} D^\dagger_{i,i} P_{i,k}$ is symmetric in $j$ and $k$ and nonnegative. Its columns sum to one: $\sum_k M_{k,j} = \sum_i P_{i,j} D^\dagger_{i,i} D_{i,i} = \sum_i P_{i,j} = 1$, because $P_{i,j} > 0$ implies $D_{i,i} > 0$. A nonnegative matrix whose columns sum to one has $\|M\|_1 = 1$, so its eigenvalues are at most $1$ in magnitude; with $B = (D^\dagger)^{1/2} P$, $M = B^\top B$ is positive semidefinite, so they are nonnegative. $M_{j,j} = \sum_i P_{i,j}^2 / D_{i,i} = 1 - O_j$ by Definition~\ref{def:ProvenanceOverlap}. For $N$, $\sum_{i'} N_{i',i} = \sum_j P_{i,j} D^\dagger_{i,i} \sum_{i'} P_{i',j} = \sum_j P_{i,j} / D_{i,i} = 1$ when $D_{i,i} > 0$. Finally, $\sum_j P_{i,j}^2 \le \sum_j P_{i,j} = D_{i,i}$ because $P_{i,j} \le 1$, with equality if and only if every nonzero $P_{i,j}$ equals $1$, i.e., every column that $\bar{z}_i$ touches is one-hot on $\bar{z}_i$. \hfill$\square$
}

\theorem[\textit{Round-Trip Bounds}]{
\label{theorem:RoundTripBounds}
For every feature-space explanation $a$ and every problem-space explanation $b$ (Definition~\ref{def:ProblemSpaceExplanation}),
\begin{equation}
    \delta_{\psi_\tau} \le 2 \sum_j |a_j| \, O_j \qquad
    \delta_{\tau_\psi} \le 2 \sum_i |b_i| \, (1 - N_{i,i})
\end{equation}
with equality when $a$ or $b$ has a single nonzero entry.
\textit{Proof.} Column $j$ of $I - M$ has entries $1 - M_{j,j}$ on the diagonal and $-M_{k,j}$ elsewhere, so by Proposition~\ref{prop:RoundTrips} its $\ell_1$ norm is $(1 - M_{j,j}) + \sum_{k \neq j} M_{k,j} = 2 (1 - M_{j,j}) = 2 O_j$. Then $\|a - Ma\|_1 = \|\sum_j a_j (I - M) e_j\|_1 \le \sum_j |a_j| \, 2 O_j$ by the triangle inequality, with equality for $a = a_j e_j$. The same argument applies to $N$ over the elements with $D_{i,i} > 0$, the only ones on which $b$ can be nonzero. \hfill$\square$
}

The bounds turn the statistics into guarantees: from $P$ alone, before any explanation is computed, they say how much of any explanation a round trip can lose, feature by feature and element by element. Their all-or-nothing cases characterize the provenance matrices that admit perfect cyclic transfers.

\theorem[\textit{Perfect Pull--Push Reconstruction}]
{
\label{theorem:PerfectPullPush}
The pull--push reconstruction of explanations of the attributable features ($a_j = 0$ for $j \in U$) is perfect ($\delta_{\psi_\tau}=0$ for every such $a$) if and only if provenance is nonoverlapping ($O=0$), i.e., no element influences two attributable features.
\textit{Proof.} ($\Leftarrow$) $O = 0$ means $O_j = 0$ for every $j \notin U$, since each $O_j \ge 0$, and Theorem~\ref{theorem:RoundTripBounds} gives $\delta_{\psi_\tau} = 0$. ($\Rightarrow$) For $a = e_j$ with $j \notin U$, Theorem~\ref{theorem:RoundTripBounds} holds with equality, so $0 = \delta_{\psi_\tau} = 2 O_j$. By Definition~\ref{def:ProvenanceOverlap}, $O_j = 0$ if and only if no element of $A_j$ influences another feature. \hfill$\square$
}

\theorem[\textit{Perfect Push--Pull Reconstruction}]
{
\label{theorem:PerfectPushPull}
The push--pull reconstruction is perfect ($\delta_{\tau_\psi} = 0$ for every problem-space explanation $b$) if and only if provenance is unambiguous ($H=0$), i.e., every column of $P$ is one-hot.
\textit{Proof.} By Theorem~\ref{theorem:RoundTripBounds}, with equality for $b = e_i$, the push--pull reconstruction is perfect if and only if $N_{i,i} = 1$ for every element with $D_{i,i} > 0$. By Proposition~\ref{prop:RoundTrips}, this holds if and only if every column of $P$ is one-hot, since every column has its support on such elements. A column is one-hot if and only if its entropy vanishes, and the unattributable columns are one-hot on the null element, so this is $H_j = 0$ for every $j$, i.e., $H = 0$. \hfill$\square$
}

While the pull--push and push--pull operations may not be practically useful for most applications, they provide an intrinsic metric we can use to evaluate the quality of explanation transfers via their respective errors $\delta_{\psi_\tau}$ and $\delta_{\tau_\psi}$, and Theorem~\ref{theorem:RoundTripBounds} ties these errors to the statistics of $P$. This metric, however, is not without limitation, as it does not permit us to determine the quality of the pullback or the pushforward in isolation. For example, if $P$'s pushforward operator is poor, $\delta_{\psi_\tau}$ could be high even if $P$'s pullback operation is strong.
